% Generated draft: requires mathematical and source review.
\documentclass[11pt]{article}
\usepackage[margin=25mm]{geometry}
\usepackage{fontspec}
\setmainfont{Latin Modern Roman}
\usepackage{amsmath,amssymb,mathtools,unicode-math}
\setmathfont{Latin Modern Math}
\usepackage{enumitem,needspace,xurl,hyperref}
\setlength{\parindent}{0pt}
\setlength{\parskip}{5pt}
\setlength{\emergencystretch}{3em}
\newcommand{\N}{\mathbb{N}}
\newcommand{\Z}{\mathbb{Z}}
\newcommand{\Q}{\mathbb{Q}}
\newcommand{\R}{\mathbb{R}}
\newcommand{\C}{\mathbb{C}}
\newcommand{\F}{\mathbb{F}}
\newcommand{\A}{\mathbb{A}}
\newcommand{\PP}{\mathbb{P}}
\newcommand{\E}{\mathbb{E}}
\newcommand{\eps}{\varepsilon}
\newcommand{\dd}{\,\mathrm{d}}
\begin{document}
\section*{UnsolvedMath ID: 30000778}
\textbf{Source problem code:} catalogue entry
\subsection{Definitions and mathematical statement}
Retain the Iwasawa-theoretic data \(\Lambda,X_K^-,Y^-,I,\psi^-,\Psi_K,\zeta,\mathcal L^-\) of the source's Conjecture 2. The requested construction is a natural \(\Lambda\)-module isomorphism
\[
 \varphi:X_K^-\xrightarrow{\ \sim\ }Y^-/IY^-.
\]
The cleaned statement explicitly uses \(X_K^-\), whereas the original abbreviated description writes \(X_K\); the restored minus-source is retained, with that source-version distinction recorded. The superscript minus denotes the source's specified eigenspace/convention and is not an invitation to replace all modules by unsigned ones.

With \(\Xi=\varphi\otimes\psi^-\), the required reciprocity compatibility is exactly
\[
 \Xi\bigl(\Psi_K(1-\zeta)^-\bigr)=\mathcal L^-
 \pmod{IY^-\otimes\Lambda^-}.
\]
Here \(1-\zeta\) is described in the original as a universal cyclotomic norm sequence and \(\mathcal L^-\) as the relevant \(p\)-adic \(L\)-function. It is not silently replaced by a single root-of-unity value or a differently normalized analytic function. The isomorphism and this specific identity are both required; equality of module ranks, a map with finite kernel and cokernel, or an unspecified scalar multiple of the compatibility is not the printed task.

The second request tests the candidate
\[
 \sigma\longmapsto\overline b(\sigma)(\alpha)
\]
obtained from the Galois action and twisted Poincar\'e duality. Its descent to the stated source, well-defined quotient target, \(\Lambda\)-linearity, isomorphism property and the displayed compatibility must all be established in the intended setup; none follows merely from naming the candidate. The class \(\alpha\), cocycle or map \(\overline b\), Galois domain and identifications are not supplied.

The arithmetic field/tower, prime, modules, ideal \(I\), minus conventions, map \(\psi^-\), reciprocity operator and tensor base/completion are missing. The congruence also needs its ambient module: if \(\varphi\) already lands in a quotient, one must know the source's lift or quotient identification to interpret reduction modulo \(IY^-\otimes\Lambda^-\). No choice of lifts, Tate twists or integral-versus-rational coefficients is invented. In particular an Iwasawa algebra or minus direct-sum decomposition is not fully determined by the symbols alone.

The literature triage checked 2026-08-21 identifies this as Sharifi's equivalent conjecture and reports constructions and verification in many cases, while retaining the full compatibility and isomorphism as open. Those cases and hypotheses are not listed. The partially-solved label and direct-source restoration note remain dated metadata, not independently verified current results.
\subsection{Short English statement}
Construct the stated integral Iwasawa-module isomorphism with its exact reciprocity-to-\(p\)-adic-\(L\)-function identity, and test the proposed Galois/duality candidate. The modules, maps, twists and tensor/congruence conventions are missing. An abstract rank match or pseudo-isomorphism would not meet the stated requirements.
\subsection{Sources}
\begin{itemize}
\item Local UnsolvedMath registry, permanent ID 30000778, source code catalogue entry, ``Iwasawa Reciprocity in Sharifi's Conjecture,'' including the original Conjecture 2 request, restored minus-module source, exact compatibility and candidate formula.
\item \emph{Algebraische Zahlentheorie (2008)}, Oberwolfach Reports, pp.\ 1741--1792, DOI 10.4171/owr/2007/30: \url{https://doi.org/10.4171/owr/2007/30}. The citation uses 2008, while the DOI and proposed-year metadata use 2007.
\item Supplied reference ``Fukaya--Kato--Sharifi, integrality of zeta elements'': \url{https://arxiv.org/abs/1507.00495}.
\item Supplied reference ``Sharifi, extension of the Fukaya--Kato method'': \url{https://doi.org/10.4171/JEMS/1519}.
\item Literature triage checked 2026-08-21; record updated 2026-09-16. No external source has been visited or independently checked here.
\end{itemize}
\end{document}
